\begin{table}[htp]
\centering
\caption{Explanation stability across ESCS quintile subgroups. Pairwise Spearman's $\rho$ rank correlation of SHAP-based feature importance rankings between ESCS quintiles. All correlations are significant at $p < 0.001$ except Q1--Q5 ($p = 3 \times 10^{-6}$).}
\label{tab:explanation-stability}
\begin{tabular}{lcc}
\toprule
Subgroup pair & Spearman's $\rho$ & $p$-value \\
\midrule
\multicolumn{3}{l}{\textbf{ESCS quintile comparisons}} \\
Q1 (lowest) vs.\ Q2 & 0.897 & $<0.001$ \\
Q1 (lowest) vs.\ Q3 & 0.851 & $<0.001$ \\
Q1 (lowest) vs.\ Q4 & 0.811 & $<0.001$ \\
Q1 (lowest) vs.\ Q5 (highest) & 0.715 & $3\times10^{-6}$ \\
Q2 vs.\ Q3 & 0.987 & $<0.001$ \\
Q2 vs.\ Q4 & 0.962 & $<0.001$ \\
Q2 vs.\ Q5 (highest) & 0.847 & $<0.001$ \\
Q3 vs.\ Q4 & 0.984 & $<0.001$ \\
Q3 vs.\ Q5 (highest) & 0.860 & $<0.001$ \\
Q4 vs.\ Q5 (highest) & 0.873 & $<0.001$ \\
\bottomrule
\end{tabular}

\par\small Feature importance rankings were computed using SHAP on the best classification model (tuned XGBoost) with 33 shared features. Higher ESCS quintiles indicate greater socioeconomic advantage. The lowest cross-group consistency (Q1 vs.\ Q5, $\rho = 0.72$) suggests modestly divergent feature importance patterns between the most and least advantaged students.
\end{table}
